\documentclass[a4paper,9pt]{article}
\usepackage[utf8]{inputenc}
\usepackage[T1]{fontenc}
\usepackage[french]{babel}
\usepackage[margin=1cm]{geometry}
\usepackage{xcolor}
\usepackage{tcolorbox}
\usepackage{enumitem}
\usepackage{multicol}
\usepackage{lmodern}
\usepackage{tabularx}
\usepackage{booktabs}
\usepackage{amsmath}

\renewcommand{\familydefault}{\sfdefault}
\pagestyle{empty}
\setlist{nosep, leftmargin=1.2em}

% --- Palette ---
\definecolor{navy}{RGB}{0,51,102}
\definecolor{teal}{RGB}{0,128,128}
\definecolor{crimson}{RGB}{180,0,0}
\definecolor{lightbg}{RGB}{245,248,252}
\definecolor{green}{RGB}{0,120,50}

\tcbset{
  basebox/.style={colback=lightbg, colframe=navy, coltitle=white,
    fonttitle=\bfseries\small, arc=2mm, boxrule=0.5mm,
    left=2mm, right=2mm, top=1mm, bottom=1mm, after skip=3pt}
}

\newcommand{\rubrique}[2]{%
  \begin{tcolorbox}[basebox, title=#1, colframe=#2]
}
\newcommand{\finrubrique}{\end{tcolorbox}}

\newcommand{\key}[1]{\textbf{\textcolor{navy}{#1}}}
\newcommand{\warn}[1]{\textbf{\textcolor{crimson}{#1}}}
\newcommand{\ok}[1]{\textbf{\textcolor{green}{#1}}}

\begin{document}

%--- TITRE ---
\begin{center}
  \colorbox{navy}{\parbox{\dimexpr\textwidth-2\fboxsep}{
    \centering\vspace{2pt}
    {\huge\bfseries\textcolor{white}{FICHE GLOBALE DE RÉVISION — BTS SIO}} \\
    {\large\textcolor{white}{Maël Mignard — BTS SIO 1\textsuperscript{ère} année — \today}}
    \vspace{2pt}
  }}
\end{center}
\vspace{4pt}

\begin{multicols}{3}

% =============================================
\begin{tcolorbox}[basebox, colframe=navy, title=RÉSEAUX — Modèle OSI]
\begin{tabular}{lll}
\textbf{\#} & \textbf{Couche} & \textbf{PDU / Proto} \\
\hline
7 & Application & HTTP, FTP, DNS \\
6 & Présentation & SSL/TLS \\
5 & Session & NetBIOS \\
4 & Transport & TCP/UDP \\
3 & Réseau & IP, ICMP \\
2 & Liaison & Ethernet, MAC \\
1 & Physique & Câbles, WiFi \\
\end{tabular}
\vspace{2pt}\\
\key{TCP} = fiable, connexion ; \key{UDP} = rapide, sans connexion \\
\key{Encapsulation} : données $\to$ segment $\to$ paquet $\to$ trame $\to$ bits
\end{tcolorbox}

% =============================================
\begin{tcolorbox}[basebox, colframe=teal, title=RÉSEAUX — Adressage IP]
\key{IPv4} : 32 bits, 4 octets (ex: 192.168.1.1) \\
\key{IPv6} : 128 bits, 8 groupes hex \\[2pt]
\textbf{Classes privées :}
\begin{itemize}
  \item 10.0.0.0/8
  \item 172.16.0.0/12
  \item 192.168.0.0/16
\end{itemize}
\key{CIDR} : /24 $\to$ 256 adresses (254 hôtes) \\
\key{Masque} : /24 = 255.255.255.0 \\
\key{Passerelle} = routeur par défaut \\
\key{DHCP} : attribution IP automatique \\
\key{DNS} : nom $\leftrightarrow$ IP
\end{tcolorbox}

% =============================================
\begin{tcolorbox}[basebox, colframe=navy, title=ROUTAGE \& VLAN]
\key{Routeur} : interconnecte réseaux (couche 3) \\
\key{Switch} : commute trames (couche 2) \\
\key{VLAN} : segmentation logique du réseau \\
\quad $\to$ VLAN 10 : Compta, VLAN 20 : IT \\
\key{Trunk} : lien inter-switch (multi-VLAN) \\
\key{Spanning Tree (STP)} : évite les boucles \\
\key{Multicast} : 1 émetteur $\to$ groupe \\[2pt]
\textbf{Commandes Cisco :}
\begin{itemize}
  \item \texttt{enable} / \texttt{conf t}
  \item \texttt{vlan 10} / \texttt{name IT}
  \item \texttt{ip address x.x.x.x mask}
  \item \texttt{show ip route}
\end{itemize}
\end{tcolorbox}

% =============================================
\begin{tcolorbox}[basebox, colframe=crimson, title=CYBERSÉCURITÉ — Piliers]
\key{CIA} (Triade) :
\begin{itemize}
  \item \textbf{C}onfidentialité : accès autorisés seulement
  \item \textbf{I}ntégrité : données non altérées
  \item \textbf{D}isponibilité : service accessible
\end{itemize}
\key{Menaces courantes :}
\begin{itemize}
  \item \warn{Phishing} : usurpation par e-mail
  \item \warn{Ransomware} : chiffrement + rançon
  \item \warn{XSS} : injection JS dans pages web
  \item \warn{SQLi} : injection SQL
  \item \warn{RCE} : exécution de code à distance
\end{itemize}
\key{RBAC} : droits selon le rôle \\
\key{MFA} : authentification multi-facteur \\
\key{DPO} : responsable RGPD
\end{tcolorbox}

% =============================================
\begin{tcolorbox}[basebox, colframe=crimson, title=CYBERSÉCURITÉ — RGPD \& PCA]
\key{RGPD} (2018) : protection données personnelles
\begin{itemize}
  \item Consentement explicite
  \item Droit à l'oubli, droit d'accès
  \item Notification breach : 72h
  \item Amende : jusqu'à 4\% CA mondial
\end{itemize}
\key{PCA} (Plan de Continuité) : maintenir activité \\
\key{PRA} (Plan de Reprise) : après sinistre \\
\key{RTO} : délai max de reprise \\
\key{RPO} : perte de données max acceptable
\end{tcolorbox}

% =============================================
\begin{tcolorbox}[basebox, colframe=green, title=ADMIN LINUX — Commandes clés]
\begin{tabular}{ll}
\texttt{ls -lah} & liste (droits, taille) \\
\texttt{chmod 755} & permissions \\
\texttt{chown u:g} & propriétaire \\
\texttt{sudo / su} & élévation droits \\
\texttt{useradd} & créer utilisateur \\
\texttt{passwd} & changer mdp \\
\texttt{ps / top} & processus \\
\texttt{grep / find} & recherche \\
\texttt{scp / ssh} & transfert sécurisé \\
\texttt{apt install} & installer paquet \\
\end{tabular}
\end{tcolorbox}

% =============================================
\begin{tcolorbox}[basebox, colframe=green, title=ADMIN LINUX — Permissions]
Format : \texttt{-rwxrwxrwx} \\
\begin{tabular}{lll}
\textbf{Valeur} & \textbf{Lettre} & \textbf{Droit} \\
4 & r & lecture \\
2 & w & écriture \\
1 & x & exécution \\
\end{tabular}
\vspace{2pt} \\
\texttt{chmod 755} = \texttt{rwxr-xr-x} \\
\warn{Jamais w+x ensemble} (risque injection) \\
\key{u}=user, \key{g}=group, \key{o}=others \\
Ex: \texttt{chmod u+x fichier.sh}
\end{tcolorbox}

% =============================================
\begin{tcolorbox}[basebox, colframe=teal, title=ALGORITHMIQUE — Structure]
\begin{verbatim}
ALGORITHME Nom
Variables
  x : Numérique
  nom : Chaîne
DEBUT
  Lire x
  Si x > 0 Alors
    Ecrire "Positif"
  Sinon
    Ecrire "Negatif"
  FinSi
FIN
\end{verbatim}
\key{Structures} : Si/Sinon, TantQue, Pour \\
\key{Types} : Numérique, Chaîne, Booléen
\end{tcolorbox}

% =============================================
\begin{tcolorbox}[basebox, colframe=teal, title=C\# — Bases]
\begin{verbatim}
int age = 18;
string nom = "Maël";
bool majeur = age >= 18;

if (majeur) {
  Console.WriteLine("Majeur");
}

for (int i = 0; i < 10; i++) {
  Console.WriteLine(i);
}
\end{verbatim}
\key{POO} : classe, objet, méthode, propriété \\
\key{Héritage} : \texttt{class B : A} \\
\key{Interface} : contrat de méthodes
\end{tcolorbox}

% =============================================
\begin{tcolorbox}[basebox, colframe=navy, title=CRÉATION WEB — HTML/CSS]
\key{HTML} : structure (balises) \\
\texttt{<header>}, \texttt{<nav>}, \texttt{<main>}, \texttt{<footer>} \\[2pt]
\key{CSS} : mise en forme \\
\texttt{display: flex / grid} \\
\texttt{position: relative / absolute} \\
\texttt{@media (max-width: 768px)} \\[2pt]
\key{Sélecteurs} : \texttt{.classe}, \texttt{\#id}, \texttt{balise} \\
\key{Box model} : margin $\to$ border $\to$ padding $\to$ content
\end{tcolorbox}

% =============================================
\begin{tcolorbox}[basebox, colframe=navy, title=CRÉATION WEB — JS \& PHP]
\key{JavaScript} (côté client) :\\
\texttt{document.getElementById("id")} \\
\texttt{addEventListener("click", fn)} \\
\texttt{fetch(url).then(r => r.json())} \\[2pt]
\key{PHP} (côté serveur) :\\
\texttt{\$\_POST["champ"]} / \texttt{\$\_GET["param"]} \\
\texttt{mysqli\_connect(host, user, pw, db)} \\
\texttt{htmlspecialchars()} : anti-XSS \\
\warn{Toujours valider/échapper les entrées}
\end{tcolorbox}

% =============================================
\begin{tcolorbox}[basebox, colframe=crimson, title=GESTION DES INCIDENTS — ITIL]
\key{Incident} : interruption non planifiée \\
\key{Problème} : cause inconnue récurrente \\
\key{Demande} : requête de service standard \\
\key{Changement} : modification planifiée \\[2pt]
\textbf{Niveaux d'impact :}
\begin{itemize}
  \item Critique : arrêt total activité
  \item Élevé : service majeur impacté
  \item Moyen : contournement possible
  \item Faible : utilisateur individuel
\end{itemize}
\key{GLPI} : outil de ticketing open-source \\
\key{SLA} : accord de niveau de service
\end{tcolorbox}

% =============================================
\begin{tcolorbox}[basebox, colframe=green, title=GESTION PATRIMOINE — GLPI]
\key{Inventaire} : parc matériel + logiciel \\
\key{Ticket} : signalement d'incident ou demande \\
\key{CMDB} : base de données de configuration \\[2pt]
\textbf{Cycle de vie matériel :}
\begin{itemize}
  \item Achat $\to$ Déploiement $\to$ Maintenance $\to$ Fin de vie
\end{itemize}
\key{Agent GLPI} : remontée inventaire auto \\
\key{Fusioninventory} : plugin d'inventaire \\
\key{Dashboard} : vue globale tickets/parc
\end{tcolorbox}

% =============================================
\begin{tcolorbox}[basebox, colframe=navy, title=MATHS — Algèbre Booléenne]
\begin{tabular}{ll}
\key{ET (AND)} & $A \cdot B$ : 1 si A=1 et B=1 \\
\key{OU (OR)} & $A + B$ : 1 si A=1 ou B=1 \\
\key{NON (NOT)} & $\overline{A}$ : inverse \\
\key{XOR} & $A \oplus B$ : différents \\
\end{tabular}
\vspace{2pt} \\
\textbf{Lois :} $A \cdot 0=0$ ; $A+1=1$ ; $\overline{\overline{A}}=A$ \\
\textbf{De Morgan :} $\overline{A \cdot B} = \overline{A}+\overline{B}$
\end{tcolorbox}

% =============================================
\begin{tcolorbox}[basebox, colframe=navy, title=MATHS — Bases de numération]
\begin{tabular}{lll}
\textbf{Base} & \textbf{Chiffres} & \textbf{Nom} \\
2 & 0,1 & Binaire \\
8 & 0-7 & Octal \\
10 & 0-9 & Décimal \\
16 & 0-9,A-F & Hexadécimal \\
\end{tabular}
\vspace{2pt} \\
\key{Décimal $\to$ Binaire} : divisions par 2 \\
\key{Binaire $\to$ Hex} : groupes de 4 bits \\
Ex: $42_{10} = 101010_2 = 2A_{16}$
\end{tcolorbox}

% =============================================
\begin{tcolorbox}[basebox, colframe=teal, title=CEJM — Économie]
\key{Agents éco} : Ménages, Entreprises, État, Banques \\
\key{Offre/Demande} : Offre$>$Demande $\to$ Prix$\downarrow$ \\
\key{Marché} : concurrentiel, oligopole, monopole \\
\key{PIB} : richesse produite par un pays \\
\key{Inflation} : hausse générale des prix \\
\key{Chômage} : personne sans emploi, cherchant activement \\
\key{Mondialisation} : interdépendance économique mondiale
\end{tcolorbox}

% =============================================
\begin{tcolorbox}[basebox, colframe=teal, title=CEJM — Droit des contrats]
\key{Contrat} : accord de volontés créant des obligations \\
\textbf{Conditions de validité :}
\begin{itemize}
  \item Consentement libre et éclairé
  \item Capacité juridique
  \item Objet licite et certain
  \item Cause licite
\end{itemize}
\key{Vices} : erreur, dol, violence \\
\key{CDI} vs \key{CDD} : durée indéterminée/déterminée \\
\key{Rupture} : démission, licenciement, rupture conventionnelle
\end{tcolorbox}

% =============================================
\begin{tcolorbox}[basebox, colframe=teal, title=CEJM — Management]
\key{Finalité entreprise} : lucrative (SA, SARL) / non lucrative (asso) \\
\key{Styles management} :
\begin{itemize}
  \item Directif : ordres, peu d'autonomie
  \item Participatif : concertation, équipe
  \item Délégatif : autonomie totale
\end{itemize}
\key{Performance} : efficacité + efficience \\
\key{RSE} : Responsabilité Sociale des Entreprises \\
\key{Parties prenantes} : actionnaires, salariés, clients, société
\end{tcolorbox}

% =============================================
\begin{tcolorbox}[basebox, colframe=green, title=ACTIVITÉS TRANSVERSALES]
\key{Git} : système de gestion de versions \\
\begin{tabular}{ll}
\texttt{git init} & initialiser dépôt \\
\texttt{git add .} & indexer modifications \\
\texttt{git commit -m} & valider \\
\texttt{git push} & envoyer \\
\texttt{git pull} & récupérer \\
\texttt{git branch} & branches \\
\end{tabular}
\vspace{2pt} \\
\key{UX/UI} : expérience/interface utilisateur \\
\key{Veille techno} : surveiller l'actualité IT \\
\key{GitLab} : plateforme DevOps (CI/CD, issues)
\end{tcolorbox}

% =============================================
\begin{tcolorbox}[basebox, colframe=green, title=ANGLAIS — Vocabulaire IT]
\begin{tabular}{ll}
\textbf{EN} & \textbf{FR} \\
\hline
network & réseau \\
firewall & pare-feu \\
server & serveur \\
database & base de données \\
backup & sauvegarde \\
patch & correctif \\
stakeholder & partie prenante \\
workflow & flux de travail \\
\end{tabular}
\end{tcolorbox}

\end{multicols}

\vspace{3pt}
\begin{center}
  \colorbox{navy}{\parbox{\dimexpr\textwidth-2\fboxsep}{
    \centering\vspace{1pt}
    \small\textcolor{white}{\textbf{BTS SIO — Option SISR (Systèmes et Réseaux) / SLAM (Solutions Logicielles)} \quad|\quad
    Épreuves : E4 (CEJM), E5 (SI), E6 (PPE/AP)}
    \vspace{1pt}
  }}
\end{center}

\end{document}
